\section{Introduction}\label{sec:intro}

Localizing a mobile robot against a prior map is, on paper, a solved
problem, but in practice it is solved in fragments. A team fielding a 2D
LiDAR reaches for \texttt{nav2\_amcl}, a particle filter that tracks a
planar pose against an occupancy grid and nothing else. A team with a 3D
LiDAR reaches instead for an \texttt{hdl\_localization}-style tracker
built around 3D NDT registration. A team that must fuse inertial and
satellite measurements reaches for a \texttt{robot\_localization}-style
node~\cite{moore2014generalized}, which composes an EKF over odometry,
IMU, and GNSS but treats the map-relative LiDAR measurement as someone
else's problem. Each of these tools carries its own topic interface,
parameter vocabulary, initialization ritual, and recovery story, so a
practitioner who moves from a 2D platform to a 3D one, or who adds a
GNSS receiver, re-integrates from scratch. The ROS~2 maintainers'
localization survey~\cite{macenski2023survey} names exactly this gap
and sketches, as unbuilt future work, a modular framework of swappable
2D and 3D localization components behind a common interface
(Section~\ref{sec:related} details this and the surrounding prior
work). The cost of the gap is not merely inconvenience: the Nav2
marathon deployment~\cite{macenski2020marathon} attributes a leading
share of its autonomy-recovery triggers to loss of localization
confidence, so how a stack initializes, degrades, and recovers is a
first-class deployment concern rather than a detail.

\prismloc{} is an engineering response to that gap. It packages three
map-based LiDAR localization backends behind one ROS~2 node contract:
\texttt{laser2d}, a 2D likelihood-field Monte Carlo localizer against an
occupancy grid; \texttt{ndt3d}, a 3D NDT-MCL backend that reuses the
same particle filter with a voxel-Gaussian measurement model over a
prior point cloud, sampling the planar $(x, y, \theta)$ pose; and
\texttt{fusion3d}, a 15-state error-state Kalman filter (ESKF) that fuses
3D LiDAR, IMU, and RTK-GNSS with IMU-bias estimation. All three present the
identical interface --- a \texttt{map}$\to$\texttt{odom} transform, a pose
with $6\times6$ covariance, and \texttt{/initialpose} seeding --- so
that choosing between a 2D scan, a 3D cloud, and an inertial/GNSS-aided
configuration is a launch-time decision rather than a change of stack.
The two MCL backends share a single pure-C++17 estimator core (motion
model, particle filter, KLD-adaptive resampling, pose estimate), and the
fusion backend is a second such core; in both cases the algorithmic code
depends only on Eigen and is unit-tested with no ROS or PCL present.
When \texttt{laser2d} is started with no initial pose, it recovers one
from a single scan by deterministic branch-and-bound (BBS) correlative
scan matching over the occupancy grid, and re-seeds the particle filter
from that result. We claim no new estimator theory: likelihood-field
MCL, NDT-MCL, ESKF fusion, and BBS scan matching are all established.
The contribution is architectural --- three known estimator families
placed behind one contract, factored so their cores are deterministically
testable off-robot --- and it is delivered as usable open-source
software.

Concretely, this paper contributes:
\begin{enumerate}
  \item A \textbf{unified, pluggable localization architecture} in which
    three estimator families (2D likelihood-field MCL, 3D NDT-MCL, and
    LiDAR+IMU+RTK-GNSS ESKF fusion) present one ROS~2
    \texttt{map}$\to$\texttt{odom} contract, so the sensor suite selects
    a backend rather than a separate stack.
  \item \textbf{ROS-free estimator cores} that hold all filter
    mathematics behind an Eigen-only, PCL-free, rclcpp-free boundary,
    making every algorithmic claim exercisable by deterministic,
    seeded-RNG unit and integration tests that run without a middleware
    runtime.
  \item A \textbf{deterministic global relocalization} path for
    \texttt{laser2d}: branch-and-bound correlative scan matching solves
    the kidnapped-robot problem from one scan and re-seeds the MCL
    posterior, replacing stochastic random-injection recovery with a
    repeatable search.
  \item An \textbf{open-source release engineered for usability},
    including fail-fast configuration validation, startup watchdog
    diagnostics that name the input that has gone silent, and a
    single-source parameter reference kept next to the code that reads
    it.
\end{enumerate}

Section~\ref{sec:related} situates \prismloc{} against prior
localization and fusion work. Section~\ref{sec:architecture} describes
the shared contract and the core/ROS boundary; Section~\ref{sec:cores}
details the estimator cores and Section~\ref{sec:ros} the ROS~2 nodes.
Section~\ref{sec:validation} presents the ROS-free testing and CI
strategy that stands in for on-robot benchmarks in this
software-description paper, and Sections~\ref{sec:limitations}
and~\ref{sec:conclusion} discuss limitations and conclude.
